\documentclass[11pt]{article}

% --- 1. Packages (Essential for ArXiv quality) ---
\usepackage[utf8]{inputenc}
\usepackage[T1]{fontenc}
\usepackage{amsmath,amssymb}
\usepackage{graphicx}
\usepackage{booktabs}
\usepackage{hyperref}
\usepackage{xcolor}
\usepackage{multirow}
\usepackage[numbers]{natbib}

% --- 2. Document Metadata ---
\title{TrucoBench: Evaluating LLM Strategic Reasoning Through Truco Paulista}
\author{William Manzoli}
\date{June 2026}

% --- 3. The Document ---
\begin{document}

\maketitle

% --- ABSTRACT ---
\begin{abstract}
We introduce \textbf{TrucoBench}, the first benchmark for evaluating large language models (LLMs) on Truco Paulista---a popular Latin American card game featuring imperfect information, bluffing, and a unique nested escalation mechanic. Unlike poker-based benchmarks, Truco's escalation system (truco $\to$ seis $\to$ nove $\to$ doze) creates a distinct strategic axis where players must reason about stake-raising, opponent modeling, and fold-or-fight decisions at multiple levels. We evaluate four frontier LLMs across a series of expert diagnostic scenarios and tournaments. Our findings reveal a significant ``Knowing-Doing Gap'': while models like GPT-4o excel at psychological tactics (bluffing/defense), they frequently fail at sequential game-theoretic logic (waste of high-value cards). TrucoBench and all evaluation code are open-source at \url{https://github.com/ManzoliW/trucobench}.
\end{abstract}

% --- INTRODUCTION ---
\section{Introduction}

Large language models have demonstrated remarkable capabilities across a wide range of tasks. A growing body of work evaluates LLMs on strategic games---particularly poker---as a testbed for reasoning under uncertainty and deception. However, the focus on poker leaves a gap: many culturally significant card games with distinct strategic properties remain unexplored as LLM benchmarks.

\textbf{Truco Paulista} is one of the most popular card games in Brazil and across Latin America. While it shares imperfect information with poker, it introduces a fundamentally different strategic axis: a \emph{nested escalation mechanic} where players can repeatedly raise the stakes of each round. Each escalation forces the opponent into a three-way decision---accept, counter-raise, or fold---creating a rich space for bluffing and meta-reasoning.

We argue that Truco Paulista offers a complementary evaluation axis:
\begin{itemize}
    \item \textbf{Escalation as nested bluffing.} Unlike poker's continuous bet sizing, Truco's discrete escalation chain creates clear decision points for commitment and credibility.
    \item \textbf{Cultural diversity in AI evaluation.} Benchmarking LLMs on games from diverse cultures helps identify potential biases in strategic reasoning.
\end{itemize}

% --- RESULTS ---
\section{Results: Diagnostic Evaluation}

The diagnostic evaluation serves as a luck-independent measure of strategic ``intuition.'' We evaluated four models against a curated set of 9 scenarios covering core Truco Paulista mechanics.

\begin{table}[ht]
\centering
\caption{Diagnostic Accuracy Across Truco Strategic Categories (\%)}
\label{tab:diagnostics}
\begin{tabular}{@{}lccccc@{}}
\toprule
Model & Overall & Bluff & Defense & Logic & Escalation \\ 
\midrule
GPT-4o-mini & 71.1 & 100.0 & 40.0 & 100.0 & 100.0 \\ 
GPT-4o & 71.1 & 100.0 & 100.0 & 100.0 & 60.0 \\ 
Heuristic (Baseline) & 36.7 & 30.0 & 0.0 & 100.0 & 100.0 \\ 
Random (Baseline) & 27.8 & 30.0 & 0.0 & 100.0 & 0.0 \\ 
\bottomrule
\end{tabular}
\end{table}

The results in Table~\ref{tab:diagnostics} reveal a significant discovery: while frontier models exhibit 100\% accuracy in bluffing and defense, they struggle with high-level sequential logic. Specifically, in scenario \texttt{MATH\_001} (Zap Management), models consistently prioritized immediate victory over round security, a sub-optimal play in professional Truco.

\section{Conclusion}
TrucoBench provides a novel lens through which to evaluate LLM reasoning. Future work will expand to 4-player team dynamics and partner signaling (``sinais'').

\end{document}
